% Part: incompleteness
% Chapter: theories-computability
% Section: intepretability

\documentclass[../../../include/open-logic-section]{subfiles}

\begin{document}

\olfileid{inc}{tcp}{itp}

\olsection{যেসব তত্ত্বে $\Th{Q}$ ব্যাখ্যাযোগ্য, সেগুলি অনির্ণেয়}

এই ফলগুলিকে আরও শক্তিশালী করা যায়। অনানুষ্ঠানিকভাবে, কোনো ভাষা
$\Lang{L_1}$-এর অন্য ভাষা~$\Lang{L_2}$-তে একটি \emph{ব্যাখ্যা} বলতে
$\Lang{L_2}$-এর !!{formula}s দিয়ে $\Lang{L_1}$-এর বিচরণক্ষেত্র,
সম্বন্ধ-প্রতীক ও অপেক্ষক-প্রতীকগুলি সংজ্ঞায়িত করাকে বোঝায়। এখানে সময়
দিয়ে কাজটি না করলেও এই সংজ্ঞাকে সুনির্দিষ্ট করা যায়।

\begin{thm}
  ধরি $\Th{T}$ এমন একটি ভাষার তত্ত্ব যাতে পাটিগণিতের ভাষাকে এমনভাবে
  ব্যাখ্যা করা যায় যে $\Th{T}$ $\Th{Q}$-এর ব্যাখ্যার সঙ্গে
  সঙ্গতিপূর্ণ। তাহলে $\Th{T}$ অনির্ণেয়। $\Th{T}$ যদি $\Th{Q}$-এর
  স্বতঃসিদ্ধগুলির ব্যাখ্যা প্রমাণ করে, তবে $\Th{T}$-এর কোনো
  সঙ্গতিপূর্ণ প্রসারণ নির্ণেয় নয়।
\end{thm}

প্রমাণটি আগের উপপাদ্যের প্রমাণের সামান্য পরিবর্তনমাত্র; একটি
প্রতিদৃষ্টান্ত ব্যবহার করে $\Th{Q}$ ও $\Th{\bar Q}$-এর একটি পৃথকীকরণ
পাওয়া যেত। সেরমেলো--ফ্র্যাঙ্কেল সেটতত্ত্বের সঙ্গে পছন্দের স্বতঃসিদ্ধ
$\Th{ZFC}$-কে সাধারণ গণিতের বহুলাংশ সম্পাদনের উপযোগী একটি স্বতঃসিদ্ধমূলক
ভিত্তি হিসেবে নেওয়া যায়। $\Th{ZFC}$-তে স্বাভাবিক সংখ্যা সংজ্ঞায়িত করা
যায়, এবং এই ব্যাখ্যায় $\Th{Q}$-এর স্বতঃসিদ্ধগুলি সত্য। তাই পাই:

\begin{cor}
$\Th{ZFC}$-এর কোনো সঙ্গতিপূর্ণ নির্ণেয় প্রসারণ নেই।
\end{cor}
% BN-SRC-242 TARGET-CORRECTION-BEGIN
\par\smallskip\noindent\textnormal{\small\textbf{উৎস-সংশোধন (BN-SRC-242):}
হিমায়িত প্রথম অনুসিদ্ধান্তে ``সঙ্গতিপূর্ণ'' শর্তটি বাদ গেছে, ফলে সেটি
মিথ্যা: কোনো অসঙ্গতিপূর্ণ প্রসারণে প্রতিটি বাক্যই উপপাদ্য, তাই সেটি
নির্ণেয়। ঠিক আগের উপপাদ্যের দ্বিতীয় বাক্য অনুসারে বাংলা অনুসিদ্ধান্তে
সঙ্গতির প্রয়োজনীয় শর্তটি পুনরুদ্ধার করা হয়েছে।}
% BN-SRC-242 TARGET-CORRECTION-END

\begin{cor}
$\Th{ZFC}$-এর কোনো সম্পূর্ণ, সঙ্গতিপূর্ণ, গণনসাধ্যভাবে
!!{axiomatizable} প্রসারণ নেই।
\end{cor}

$\Th{ZFC}$-এর ভাষায় শুধু একটি দ্বি-স্থানীয় সম্বন্ধ~$\in$ আছে।
(বস্তুত, সমতাও দরকার হয় না।) তাই পাই:

\begin{cor}
দ্বি-স্থানীয় সম্বন্ধ-প্রতীকবিশিষ্ট যেকোনো ভাষার প্রথম-ক্রমের যুক্তিবিদ্যা
অনির্ণেয়।
\end{cor}

ফলটি দুইটি এক-স্থানীয় অপেক্ষক-প্রতীকবিশিষ্ট যেকোনো ভাষার ক্ষেত্রেও
প্রযোজ্য, কারণ এগুলি দিয়ে একটি দ্বি-স্থানীয় সম্বন্ধ-প্রতীককে অনুকরণ করা
যায়। এইমাত্র উল্লিখিত ফলগুলির সীমা সূক্ষ্ম: দেখা যায়, কেবল
\emph{এক-স্থানীয়} সম্বন্ধ-প্রতীক এবং সর্বাধিক একটি \emph{এক-স্থানীয়}
অপেক্ষক-প্রতীকবিশিষ্ট ভাষার প্রথম-ক্রমের যুক্তিবিদ্যা নির্ণেয়।

আরও একটি তথ্য। আমরা জানি, $\Obj 0$, $'$, $+$, $\times$, $<$ ভাষার
যেসব বাক্য মানক মডেলে সত্য, তাদের সেট অনির্ণেয়। বস্তুত, অন্য প্রতীকগুলির
সাহায্যে $<$-কে এবং পরে $\times$ ও $'$-এর সাহায্যে $+$-কে সংজ্ঞায়িত
করা যায়। তাই $\Obj 0$, $'$, $\times$ ভাষায় মানক মডেলে সত্য বাক্যগুলির
সেট অনির্ণেয়। অন্যদিকে, প্রেসবুর্গার দেখিয়েছেন যে $\Obj 0$, $'$, $+$
ভাষায় মানক মডেলে সত্য বাক্যগুলির সেট নির্ণেয়। তবে সেই পদ্ধতি গণনার
দিক থেকে ব্যবহারিক নয়।
% BN-SRC-243 TARGET-CORRECTION-BEGIN
\par\smallskip\noindent\textnormal{\small\textbf{উৎস-সংশোধন (BN-SRC-243):}
হিমায়িত শেষ অনুচ্ছেদে প্রেসবুর্গার পাটিগণিতের বাক্যগুলিকে ``পাটিগণিতের
ভাষায় সত্য'' বলা হয়েছে; কিন্তু কোনো ভাষা নিজে বাক্যের সত্যতা নির্ধারণ
করে না। অনুচ্ছেদের আগের দুই দাবির সঙ্গে মিলিয়ে বাংলা পাঠে মানক মডেলে
সত্যতার অর্থটি স্পষ্ট করা হয়েছে।}
% BN-SRC-243 TARGET-CORRECTION-END

\end{document}
